\textbf{One ChromData store serves tracing data at real scale, alongside the formats the community uses.}
\textbf{a}~The ChromData data model: \pkg{spots} (one row per observation) point at a shared locus axis (\pkg{bins}) through \pkg{bin\_id}, and at \pkg{traces} and \pkg{cells} through \pkg{trace\_id} and \pkg{cell\_id}; coordinates (+ layers) and per-spot tracks are row-aligned with the spots, and contact maps are linked by bins and a cell map rather than copied.
Right: the \pkg{x.chromdata.zarr} layout (format 2.2) as written by \pkg{uchrom.core.zarrcd}; the glyph at each node gives its storage type (Parquet table, Zarr array, JSON attributes).
Every spot-aligned value is stored once: coordinates and keys in one Parquet table per chromosome, every other spot-aligned column in cell-sorted primary tables, and spot columns that only repeat a value of their bin, cell or trace once per key; \pkg{index/} maps the (cell, trace, chromosome) runs of the two sides one to one.
\textbf{b}~Interoperability with community formats: what survives an export followed by an import (values compared exactly, coordinates to the format's precision). Filled, preserved; half-filled, partial, or not representable in that format; open, not applicable. PDB is export-only and .3dg import-only in U-Chrom.
\textbf{c}~File size and full-read time (into memory) vs number of spots for whole-cell subsets of the Takei et al.\ 2025 cerebellum DNA seqFISH+ data (62 per-spot tracks)~\citep{takei2025}, for chromdata.zarr 2.2 and the storage schemes publicly used by the community for chromatin-tracing data: the 4DN FOF-CT core table (CSV), AnnData with spots as observations, and SpatialData (spots as a 3-D points element, cells as a table; \pkg{spatialdata} 0.7.2 defaults, full read = \pkg{read\_zarr} + \pkg{.compute()}); median of 3 fresh processes.
At $10^7$ spots chromdata.zarr stores 1.8\,GB and reads in 0.57\,s (FOF-CT 6.7\,GB, 99\,s; AnnData 5.6\,GB, 5.4\,s; SpatialData 2.2\,GB, 2.2\,s).
Flat Parquet, the storage engine under chromdata.zarr and therefore its lower bound without hierarchy, index or metadata, and the ChromData format versions are in Supplementary Fig.~S1.
\textbf{d}~Time per query vs spots: backed chromdata.zarr from disk (\pkg{get\_trace}, \pkg{get\_cell}, \pkg{get\_chrom} with all 62 tracks and with coordinates only, and opening the store; median of 20 random ids $\times$ 3 fresh processes, shading IQR) vs loading the whole table in memory, then subsetting (FOF-CT, AnnData; load + one \pkg{get\_trace}) and SpatialData's lazy route (open + a dask filter on \pkg{trace\_id} + \pkg{.compute()}).
At $10^7$ spots \pkg{get\_trace} takes 6.2\,ms and \pkg{get\_cell} 9.0\,ms; \pkg{get\_chrom} with all tracks reads a slice of every cell-sorted row group (0.52\,s; 13\,ms with coordinates only).
\textbf{e}~Beyond RAM. Left: memory used to open the data and answer the trace and cell queries of \textbf{d} (above the interpreter's baseline), backed chromdata.zarr vs the whole table loaded (colours as in \textbf{d}).
Right: exact median 3-D distance map of chr11, peak memory vs traces when the store is loaded whole and when \pkg{iter\_traces} streams it under a 2\,GB budget (1.55 vs 12.3\,GB at $10^7$ spots); streamed and in-memory medians are bitwise identical at every size (inset: $10^7$ spots, 10.7\,M bin pairs; at $10^8$ against the chromosome loaded alone).
Open markers ($10^8$ spots in \textbf{d} and \textbf{e}): the Takei data replicated ten times with renamed cells and traces, not new cells; at $10^8$ loading the whole table does not fit in 64\,GB of RAM.
\textbf{f}~Groups from one modality, structure from another (scHiCAR mouse primary motor cortex, 5,313 cells)~\citep{schicar}.
Leiden clustering computed in ChromData on the RNA PCA (16 clusters; not the published labels) groups the cells; three contrasting clusters are shown on the RNA UMAP: oligodendrocytes (\textit{Plp1}, \textit{St18}, \textit{Prr5l}, \textit{Mobp}; $n=294$), astrocytes (\textit{Gpc5}, \textit{Slc1a3}, \textit{Atp1a2}; $n=448$) and neurons (no glial or vascular marker expression; $n=1{,}336$).
Their pseudo-bulk contact maps are summed from the linked per-cell .scool through the shared index (the path the web browser uses; chr2:40--90\,Mb, 1\,Mb bins, each normalised to its cis contacts).
Contact probability vs distance, $P(s)$, over all chromosomes: slopes over 2--30\,Mb of $-0.83$ (neurons), $-0.67$ (oligodendrocytes) and $-0.67$ (astrocytes), a modest difference.
Pseudo-bulk from the per-cell .scool equals the published per-type .cool for 22/22 cell types (max $|\Delta| = 0$ over 4.2\,M pixels, all chromosomes).
Right: time to sum the pseudo-bulk of chr1 vs cells aggregated, first query and cached.
Machine: Apple M5 Pro (18 cores), 64\,GB RAM; warm OS page cache.
Zarr and the Zarr logo are marks of the Zarr project (NumFOCUS); Apache Parquet and the Parquet logo are trademarks of The Apache Software Foundation; all logos are used for identification only and do not imply endorsement.
